\begin{examplebox}
\textbf{例1：椭圆中直接核对方向。}
$x^2/4+y^2/3=1$，$P=(1,3/2)$，焦点为 $(\pm1,0)$。
$PF_1=5/2$，$PF_2=3/2$，两个焦点方向单位向量为
$u=(-4/5,-3/5)$、$v=(0,-1)$。
$u+v=(-4/5,-8/5)$ 的斜率为 $2$，恰为法线斜率；切线 $x+2y=4$ 的斜率为 $-1/2$。
\par\medskip\noindent
\textbf{例2：双曲线中内角平分线为切线。}
$x^2-y^2/3=1$，$P=(2,3)$，焦点为 $(\pm2,0)$。
$u=(-4/5,-3/5)$、$v=(0,-1)$，其和的斜率为 $2$。
切线 $2x-y=1$ 的斜率正是 $2$，所以切线平分内部焦角。
\par\medskip\noindent
\textbf{例3：抛物线反射路径。}
$y^2=4x$，$F=(1,0)$，$P=(1,2)$，$A=(5,2)$。
$PF$ 向下，$PA$ 向右，内部角为 $90^\circ$，法线方向 $(1,-1)$ 平分该角。
切线 $y=x+1$ 是外角平分线，$FP+PA=2+4=6$。
把 $F$ 关于切线反射到 $F'=(-1,2)$，则 $F',P,A$ 共线，且
$F'P+PA=F'A=6$，与反射定律一致。
\end{examplebox}
